\documentclass[11pt,a4paper]{article}
\usepackage[margin=2.4cm]{geometry}
\usepackage{amsmath,amssymb}
\usepackage{booktabs}
\usepackage{parskip}
\usepackage[hidelinks]{hyperref}

\newcommand{\cd}{c^\dagger}
\newcommand{\ad}{a^\dagger}
\newcommand{\up}{\uparrow}
\newcommand{\dn}{\downarrow}

\begin{document}

\begin{center}
{\large\bfseries VQE for the half-filled Hubbard ring on MIMIQ:\\ procedure and first results}\\[6pt]
Khuram Shahzad, October 2026
\end{center}

%-----------------------------------------------------------------------------
\section{Goal}

Prepare the ground state of the Hubbard ring with the
variational quantum eigensolver (VQE)~\cite{peruzzo} on the MIMIQ simulator~\cite{mimiq}, then compute its
$n$BRDMs and GME and compare them with exact diagonalization (ED).
System: $L = 6$ sites, $K = 6$ electrons (half filling), $S_z = 0$, $V = 0$,
$t = 1$; this needs 12 qubits. Code, tests and results:
\url{https://github.com/ItsKhuramShahzad/mimiq-gme}.

\begin{center}
{\renewcommand{\arraystretch}{1.6}
\begin{tabular}{ll}
  \toprule
  Symbol & Meaning \\
  \midrule
  $E_0^{\rm ED}$,\ \ $|\Psi_0^{\rm ED}\rangle$ & exact ground-state energy and state, from exact diagonalization (ED) \\
  $E_1^{\rm ED}$ & exact energy of the first excited state (ED) \\
  $\Delta = E_1^{\rm ED} - E_0^{\rm ED}$ & energy gap \\
  $E_{\rm VQE}$,\ \ $|\Psi(\boldsymbol\theta^*)\rangle$ & lowest energy found by VQE and its state \\
  $E_{\rm HF}$ & energy of the Hartree--Fock starting state \\
  \bottomrule
\end{tabular}}
\end{center}

%-----------------------------------------------------------------------------
\section{Model}

\begin{gather}
H_e = - t \sum_{k=1}^L \sum_{\sigma} \big(\cd_{k,\sigma} c_{k+1,\sigma} + \cd_{k+1,\sigma} c_{k,\sigma}\big)
+ U \sum_{k=1}^L n_{k,\up}\, n_{k,\dn} \nonumber\\
+ V \sum_{k=1}^L n_k\, n_{k+1} + \epsilon \sum_{k=1}^L n_k , \label{eq:H}
\end{gather}
with $n_{k,\sigma} = \cd_{k,\sigma}c_{k,\sigma}$, $n_k = n_{k,\up} + n_{k,\dn}$
and site $L+1 \equiv 1$ (ring). At fixed $K$ the $\epsilon$ term only adds the
constant $\epsilon K$, so $\epsilon = 0$ is used. Since
$H_e(t,U,V) = t\,H_e(1,U/t,V/t)$, the ground state depends only on $U/t$ and $V/t$.

This Hamiltonian is never changed. VQE does not diagonalize it; it only
measures its expectation value.

%-----------------------------------------------------------------------------
\section{From electrons to qubits (Jordan--Wigner)}

One qubit per spin orbital; $|1\rangle$ = occupied, $|0\rangle$ = empty:
\begin{equation}
  q(k,\sigma) = 2(k-1) + \sigma , \qquad \sigma = 0\ (\up),\ 1\ (\dn),
\end{equation}
so the 12 qubits are $1\up, 1\dn, 2\up, 2\dn, \dots, 6\dn$. The
Jordan--Wigner transformation~\cite{jw} is
\begin{equation}
  c_{k,\sigma} = Z_0 \cdots Z_{q-1}\,\tfrac12 (X_q + iY_q) , \qquad
  n_{k,\sigma} = \tfrac12 (1 - Z_q) .
\end{equation}
The $Z$ string gives the fermion sign. The terms of Eq.~(\ref{eq:H}) become
($q < q'$)
\begin{align}
  \cd_{q} c_{q'} + \cd_{q'} c_{q} &= \tfrac12 \big( X_q Z_{q+1}\cdots Z_{q'-1} X_{q'} + Y_q Z_{q+1}\cdots Z_{q'-1} Y_{q'} \big) , \\
  n_q\, n_{q'} &= \tfrac14 \big( 1 - Z_q - Z_{q'} + Z_q Z_{q'} \big) .
\end{align}
So the Hamiltonian is a sum of Pauli strings with real coefficients,
\begin{equation}\label{eq:Hpauli}
  H_e = \lambda_0 + \sum_\mu \lambda_\mu P_\mu ,
\end{equation}
with 42 strings for $L = 6$ (24 hopping, 18 on-site) and $\lambda_0 = UL/4$.
This step is exact: the spectrum does not change.

%-----------------------------------------------------------------------------
\section{How VQE works}

A circuit with parameters $\boldsymbol\theta$ (the \emph{ansatz}) prepares a
trial state $|\Psi(\boldsymbol\theta)\rangle$. MIMIQ measures
\begin{equation}
  E(\boldsymbol\theta) = \langle \Psi(\boldsymbol\theta)| H_e |\Psi(\boldsymbol\theta)\rangle
  = \lambda_0 + \sum_\mu \lambda_\mu \langle P_\mu\rangle_{\boldsymbol\theta}
  \;\ge\; E_0^{\rm ED} .
\end{equation}
A classical optimizer (COBYLA~\cite{cobyla}, which needs no gradients) changes
$\boldsymbol\theta$ to lower $E(\boldsymbol\theta)$, and the loop repeats until
the energy stops decreasing. By the variational principle the lowest energy
found is an upper bound to $E_0^{\rm ED}$, and the state approaches $|\Psi_0^{\rm ED}\rangle$.

%-----------------------------------------------------------------------------
\section{Initial state}

The circuit starts from the ground state at $U = 0$, which is easy to prepare.

\emph{Orbitals.} Diagonalize the $L \times L$ one-electron hopping matrix
$h_{kk'} = -t(\delta_{k',k+1} + \delta_{k',k-1})$:
\begin{equation}
  h = C\,\mathrm{diag}(e_1,\dots,e_L)\,C^T , \qquad
  c_{k,\sigma} = \sum_{p} C_{kp}\, a_{p,\sigma} , \qquad
  e_p = -2t, -t, -t, t, t, 2t .
\end{equation}
The orbitals are only used to build the start; the energy is always measured
with the full $H_e$.

\emph{Hartree--Fock.} Fill the three lowest orbitals with both spins:
\begin{equation}
  |\Phi_0\rangle = \prod_{p=1}^{3} \ad_{p,\up}\ad_{p,\dn}|\mathrm{vac}\rangle
  = X_0 X_1 \cdots X_5 |0\cdots0\rangle , \qquad
  E_{\rm HF} = -8t + \tfrac32 U .
\end{equation}
This is a closed shell for $L = 6$.

\emph{Back to the sites.} The circuit $\mathcal{R}_C$ rewrites the state from
orbitals to sites,
\begin{equation}
  \mathcal{R}_C\, \ad_{p,\sigma}\, \mathcal{R}_C^\dagger = \sum_k C_{kp}\,\cd_{k,\sigma} .
\end{equation}
It is a product of 15 Givens rotations of neighboring orbitals~\cite{kivlichan}; each one is
exactly two Pauli rotations per spin, so there is no approximation. The state
$\mathcal{R}_C|\Phi_0\rangle$ is the exact $U = 0$ ground state on the sites.
Its overlap with the true ground state is $F = 0.980$, $0.689$, $0.265$ at
$U/t = 1, 4, 10$: a good start at small $U/t$, a poor one at large $U/t$.

%-----------------------------------------------------------------------------
\section{Ans\"atze}

\emph{UCCSD}~\cite{romero} (117 parameters). Single and double excitations from occupied to
empty orbitals, $|\Psi\rangle = \mathcal{R}_C\, e^{T(\boldsymbol\theta)-T(\boldsymbol\theta)^\dagger}|\Phi_0\rangle$,
applied with one first-order Trotter step.

\emph{LUCJ}~\cite{motta} (130 parameters). Local unitary cluster Jastrow: orbital rotations
and density--density phases, two repetitions.

\emph{HVA} (Hamiltonian variational ansatz~\cite{wecker,cade}, $3p$
parameters). Split $H_e$ into three groups whose terms commute: the on-site
term $H_U = \sum_k n_{k\up}n_{k\dn}$, the hopping on the bonds $(1,2)$, $(3,4)$,
$(5,6)$ ($H_A$), and the hopping on the bonds $(2,3)$, $(4,5)$, $(6,1)$ ($H_B$).
Then
\begin{equation}
  |\Psi(\boldsymbol\theta)\rangle = \prod_{l=1}^{p}
  e^{-i\theta^{(l)}_B H_B}\, e^{-i\theta^{(l)}_A H_A}\, e^{-i\theta^{(l)}_U H_U}\;
  \mathcal{R}_C|\Phi_0\rangle .
\end{equation}
Inside each group all Pauli strings commute, so every layer is exact (no
Trotter error). Every layer conserves $K$, $S_z$ and $S^2$, so the state stays a
singlet. At $\boldsymbol\theta = 0$ it is the $U = 0$ ground state.

\emph{Warm start along $U/t$.} The HVA circuit does not depend on $U$, and the
ground state changes smoothly with $U/t$ (the gap stays open). Therefore the
points $U/t = 1, 2, 3, 4, 5, 6, 8, 10$ are solved in this order, each one
starting from the optimal angles of the previous one, $\boldsymbol\theta_{\rm start}(U_{j+1}) = \boldsymbol\theta^*(U_j)$.

%-----------------------------------------------------------------------------
\section{Accuracy criterion}

Two quantities are reported, the energy error and the fidelity (the fidelity
is available on the simulator):
\begin{equation}
  \Delta E = E_{\rm VQE} - E_0^{\rm ED} , \qquad
  F = |\langle\Psi_0^{\rm ED}|\Psi(\boldsymbol\theta^*)\rangle|^2 .
\end{equation}
They are linked by the gap $\Delta = E_1^{\rm ED} - E_0^{\rm ED}$:
\begin{equation}
  F \;\ge\; 1 - \frac{\Delta E}{\Delta} .
\end{equation}
At large $U/t$ the gap is small ($\Delta = 1.66$ at $U/t = 1$, $0.27$ at
$U/t = 10$), so a small $\Delta E$ alone does not guarantee a good state. The proposed
criterion is $\Delta E < 0.01\,\Delta$ (which gives $F \ge 0.99$),
together with $|\Delta \mathcal{C}_{\rm GME}| < 0.01$.

%-----------------------------------------------------------------------------
\section{Results ($L = 6$, $t = 1$, $V = 0$)}

\begin{center}
\begin{tabular}{llrcc}
  \toprule
  $U/t$ & ansatz & parameters & $\Delta E$ & $F$ \\
  \midrule
  1  & UCCSD                         & 117 & $1.6\times10^{-4}$ & 0.99997 \\
     & LUCJ                          & 130 & $4.9\times10^{-3}$ & 0.9990 \\
     & HVA, $p = 8$                  &  24 & $2.8\times10^{-4}$ & 0.99994 \\
  \midrule
  10 & UCCSD                         & 117 & $3.9\times10^{-1}$ & 0.942 \\
     & LUCJ                          & 130 & $2.1\times10^{-1}$ & 0.634 \\
     & HVA, $p = 8$, start $\boldsymbol\theta = 0$ & 24 & $6.6\times10^{-2}$ & 0.952 \\
     & HVA, $p = 8$, warm start      &  24 & $3.8\times10^{-3}$ & \textbf{0.99937} \\
  \bottomrule
\end{tabular}
\end{center}

HVA with $p = 8$ and the warm start along $U/t$:
\begin{center}
\begin{tabular}{ccccc}
  \toprule
  $U/t$ & $\Delta$ & $\Delta E$ & $\Delta E/\Delta$ & $F$ \\
  \midrule
  1  & 1.656 & $2.4\times10^{-3}$ & 0.15\,\% & 0.99954 \\
  2  & 1.313 & $2.1\times10^{-3}$ & 0.16\,\% & 0.99957 \\
  3  & 1.009 & $3.8\times10^{-3}$ & 0.37\,\% & 0.99922 \\
  4  & 0.770 & $3.9\times10^{-3}$ & 0.51\,\% & 0.99929 \\
  5  & 0.601 & $4.1\times10^{-3}$ & 0.68\,\% & 0.99929 \\
  6  & 0.485 & $3.3\times10^{-3}$ & 0.67\,\% & 0.99942 \\
  8  & 0.349 & $3.6\times10^{-3}$ & 1.03\,\% & 0.99944 \\
  10 & 0.274 & $3.8\times10^{-3}$ & 1.39\,\% & 0.99937 \\
  \bottomrule
\end{tabular}
\end{center}

Main observations:
\begin{itemize}
  \item At small $U/t$ every ansatz works.
  \item At large $U/t$ a lower energy does not mean a better state: LUCJ has a
    lower energy than UCCSD at $U/t = 10$ but a much lower fidelity, because it
    mixes in low-lying spin excitations that cost little energy.
  \item HVA with $p = 8$ and the warm start gives $F \ge 0.9992$ from $U/t = 1$ to
    $10$ with only 24 parameters. The criterion $\Delta E < 0.01\,\Delta$ holds
    up to $U/t = 6$ and is just missed at $U/t = 8$ and $10$, where the gap is small.
  \item A test of entanglement: the entanglement entropy of 1, 2 and 3 sites of
    the HVA states agrees with ED within $0.02$ bits.
\end{itemize}

%-----------------------------------------------------------------------------
\section{Next steps}

\begin{enumerate}
  \item $n$BRDMs of the HVA states measured from Pauli strings on MIMIQ, GME
    with the GCG solver, compared with ED at every $U/t$.
  \item A cheaper circuit for the digital twin~\cite{consiglio}: ordering the qubits as
    $1\up,\dots,6\up,1\dn,\dots,6\dn$ makes every hop (except across the ring
    bond) a two-qubit gate, which halves the number of two-qubit gates of the
    $p = 8$ circuit (1616 $\to$ 856), with exactly the same state.
\end{enumerate}

\begin{thebibliography}{99}
\small
\bibitem{peruzzo} A. Peruzzo \emph{et al.}, A variational eigenvalue solver on a photonic
  quantum processor, Nat.\ Commun.\ \textbf{5}, 4213 (2014).
\bibitem{mimiq} MIMIQ quantum circuit simulator, QPerfect, \url{https://qperfect.io}.
\bibitem{jw} P. Jordan and E. Wigner, \"Uber das Paulische \"Aquivalenzverbot,
  Z.\ Phys.\ \textbf{47}, 631 (1928).
\bibitem{cobyla} M. J. D. Powell, A direct search optimization method that models the
  objective and constraint functions by linear interpolation, in \emph{Advances in
  Optimization and Numerical Analysis} (Kluwer, 1994), pp.\ 51--67.
\bibitem{kivlichan} I. D. Kivlichan \emph{et al.}, Quantum simulation of electronic
  structure with linear depth and connectivity, Phys.\ Rev.\ Lett.\ \textbf{120}, 110501 (2018).
\bibitem{romero} J. Romero \emph{et al.}, Strategies for quantum computing molecular
  energies using the unitary coupled cluster ansatz, Quantum Sci.\ Technol.\ \textbf{4},
  014008 (2019).
\bibitem{motta} M. Motta \emph{et al.}, Bridging physical intuition and hardware
  efficiency for correlated electronic states: the local unitary cluster Jastrow ansatz
  for electronic structure, Chem.\ Sci.\ \textbf{14}, 11213 (2023).
\bibitem{wecker} D. Wecker, M. B. Hastings and M. Troyer, Progress towards practical
  quantum variational algorithms, Phys.\ Rev.\ A \textbf{92}, 042303 (2015).
\bibitem{cade} C. Cade, L. Mineh, A. Montanaro and S. Stanisic, Strategies for solving
  the Fermi--Hubbard model on near-term quantum computers, Phys.\ Rev.\ B \textbf{102},
  235122 (2020).
\bibitem{consiglio} M. Consiglio \emph{et al.}, Variational quantum eigensolver for
  SU($N$) fermions, J.\ Phys.\ A \textbf{55}, 265301 (2022).
\end{thebibliography}

\end{document}
